%
% client_calls_locks_wait_held_buggy.tex
%
% FIDELITY VARIANT of docs/client_calls_locks.tex.  The caller waits for its
% reply while holding the table lock.  C6 must fail: waitHeld = fset is
% reachable.
%
% (The text below is the intact model's.)
%
% Companion to docs/client_calls.tex: the LOCK ORDER DES-107 requires of the
% Hub's request relay, checked for deadlock-freedom.  Same method as
% docs/listener_locks.tex.
%
% client_calls.tex makes each critical section one atomic step, so its
% deadlock check says nothing about locks.  The claim that the relay cannot
% deadlock is made here, and here alone.  The locks (DES-107 decision 3, C6,
% and DES-101):
%
%   sto  the store lock: the ClientLifecycle steps run under it.
%   rgy  the ClientRegistry lock: the registrations and the leases.
%   tbl  the ClientCalls table lock: the pending calls.  A leaf.
%
% THE ORDER:  sto < rgy < tbl.  A rank function over the locks is strictly
% increasing along every edge the programs take, so the acquisition graph is
% acyclic.  The model is not that argument.  It runs the four programs below,
% each as a thread of acquire / release / wait / signal steps over the locks,
% against every interleaving, and asks ProB for a state in which some thread
% is not done and none can move.
%
% MUTUAL EXCLUSION IS ASSUMED, AS A LOCK IS MODELLED: Acquire is guarded by
% "the lock is free".  What is checked:
%   (1) deadlock-freedom: Halt is enabled only when every thread is done;
%   (2) ORDER, as a negated goal: orderBroken is set by an Acquire made while
%       the thread holds a lock of equal or higher rank;
%   (3) C6, as two negated goals: waitHeld is set by a wait made while the
%       thread holds any lock, wakeHeld by a wake made while it holds the
%       table lock.
% None of the three is a conjunct of the state schema: ProB treats the primed
% state schema as part of each operation's guard, so a violation written there
% would disable the operation instead of reporting it.
%
% THE WAIT.  A caller blocks on its own event for at most the deadline.  There
% is one event per thread (the set woken): a Signal or a Deadline sets one
% thread's event, and the wait step of that thread alone is enabled by it and
% consumes it.  The Deadline step of a thread is enabled while it is at its
% wait, so a wait always ends.  What the model asks is not whether the wait
% ends but what the waiter holds while it waits: nothing.
%
% THE PROGRAMS (the index of the program counter numbers the steps):
%   callC     rgy{ find the provider; tbl{enter the call, check the limit} };
%             wait holding nothing.  The entry is made under the registry
%             lock: were the lock released between finding the provider and
%             entering the call, the provider's session could end in the gap
%             and the call would be pending on a dead session (C4).
%   resolver  tbl{match and remove the entry}; wake.  The reader on a reply or
%             a refusal frame, and equally the deadline timer.
%   endC      sto{ rgy{ withdraw the registrations; tbl{resolve the pending
%             calls lost}; wake } }.  The wake is after the table lock is
%             released, inside the store and registry locks: DES-107 C6.
%             Disconnect, lease lapse, supersession and teardown are this
%             program.
%   regC      sto{ rgy{ register } }: a service register is a store-locked
%             lifecycle step, never a bare registry write.
%
% THE REJECTED ORDERS: client_calls_locks_inverted_buggy.tex has the lifecycle
% end take the table lock before the registry lock; a caller holds the registry
% lock and wants the table, the lifecycle holds the table and wants the
% registry.  client_calls_locks_wait_held_buggy.tex has the caller wait while
% holding the table lock.  client_calls_locks_wake_held_buggy.tex has the
% lifecycle wake the waiter while holding the table lock.
%
% Type-check:  fuzz docs/client_calls_locks.tex
% Model-check: scripts/verify_client_calls_model.sh
%
\documentclass[a4paper,10pt,fleqn]{article}
\usepackage[margin=1in]{geometry}
\usepackage{fuzz}

\begin{document}

\title{lux Client Calls: the Lock Order}
\maketitle

\begin{zed}
  [T]
\end{zed}

A thread: a caller, the reader or timer that resolves a call, the lifecycle
step that ends a session, the lifecycle step that registers a service.  Two
suffice for a cycle; three give every pair a third party.

\begin{zed}
  LOCK ::= sto | rgy | tbl
\end{zed}

\begin{zed}
  ACT ::= grab | drop | wt | sg | fin
\end{zed}

\begin{zed}
  PROG ::= callC | resolver | endC | regC
\end{zed}

\begin{zed}
  FSTATE ::= fclear | fset
\end{zed}

\begin{zed}
  TST ::= tnew | trun | tdone
\end{zed}

Steps: acquire or release a lock, wait on the caller's event, set that event
(wake the waiter), finish.

\begin{axdef}
  lvl : LOCK \fun 1 \upto 3
\where
  lvl = \{ sto \mapsto 1, rgy \mapsto 2, tbl \mapsto 3 \}
\end{axdef}

The order of DES-107: store, then registry, then table.  The order is a
ranking, not a premise of the model: it is the standard the flag
\texttt{orderBroken} measures each acquisition against.

\begin{axdef}
  act : PROG \cross (1 \upto 8) \pfun ACT \\
  lck : PROG \cross (1 \upto 8) \pfun LOCK
\where
  act = \{ (callC, 1) \mapsto grab, (callC, 2) \mapsto grab,
    (callC, 3) \mapsto drop, (callC, 4) \mapsto wt, (callC, 5) \mapsto drop,
    (callC, 6) \mapsto fin,
    (resolver, 1) \mapsto grab, (resolver, 2) \mapsto drop,
    (resolver, 3) \mapsto sg, (resolver, 4) \mapsto fin,
    (endC, 1) \mapsto grab, (endC, 2) \mapsto grab, (endC, 3) \mapsto grab,
    (endC, 4) \mapsto drop, (endC, 5) \mapsto sg, (endC, 6) \mapsto drop,
    (endC, 7) \mapsto drop, (endC, 8) \mapsto fin,
    (regC, 1) \mapsto grab, (regC, 2) \mapsto grab, (regC, 3) \mapsto drop,
    (regC, 4) \mapsto drop, (regC, 5) \mapsto fin \} \\
  lck = \{ (callC, 1) \mapsto rgy, (callC, 2) \mapsto tbl,
    (callC, 3) \mapsto rgy, (callC, 5) \mapsto tbl,
    (resolver, 1) \mapsto tbl, (resolver, 2) \mapsto tbl,
    (endC, 1) \mapsto sto, (endC, 2) \mapsto rgy, (endC, 3) \mapsto tbl,
    (endC, 4) \mapsto tbl, (endC, 6) \mapsto rgy, (endC, 7) \mapsto sto,
    (regC, 1) \mapsto sto, (regC, 2) \mapsto rgy, (regC, 3) \mapsto rgy,
    (regC, 4) \mapsto sto \}
\end{axdef}

The programs, one row of \texttt{act} and \texttt{lck} per step.

\begin{schema}{LS}
  ts : T \fun TST \\
  pr : T \fun PROG \\
  pc : T \fun 1 \upto 8 \\
  owner : LOCK \pfun T \\
  woken : \power T \\
  orderBroken : FSTATE \\
  waitHeld : FSTATE \\
  wakeHeld : FSTATE
\end{schema}

\texttt{owner} is the lock table; a lock in its domain is held.
\texttt{woken} is the set of threads whose event is set: one event per
caller, as DES-107 gives each call its own outcome.

\begin{schema}{Init}
  LS'
\where
  ts' = T \cross \{ tnew \} \\
  pr' = T \cross \{ callC \} \\
  pc' = T \cross \{ 1 \} \\
  owner' = \emptyset \\
  woken' = \emptyset \\
  orderBroken' = fclear \\
  waitHeld' = fclear \\
  wakeHeld' = fclear
\end{schema}

\begin{schema}{Start}
  \Delta LS \\
  t? : T \\
  p? : PROG
\where
  ts~t? = tnew \\
  ts' = ts \oplus \{ t? \mapsto trun \} \\
  pr' = pr \oplus \{ t? \mapsto p? \} \\
  pc' = pc \\
  owner' = owner \\
  woken' = woken \\
  orderBroken' = orderBroken \\
  waitHeld' = waitHeld \\
  wakeHeld' = wakeHeld
\end{schema}

A thread picks its program once.

\begin{schema}{AcquireOrdered}
  \Delta LS \\
  t? : T
\where
  ts~t? = trun \\
  act~(pr~t?, pc~t?) = grab \\
  lck~(pr~t?, pc~t?) \notin \dom owner \\
  \forall l : \dom owner | owner~l = t? \spot
    lvl~l < lvl~(lck~(pr~t?, pc~t?)) \\
  owner' = owner \cup \{ lck~(pr~t?, pc~t?) \mapsto t? \} \\
  pc' = pc \oplus \{ t? \mapsto pc~t? + 1 \} \\
  ts' = ts \\
  pr' = pr \\
  woken' = woken \\
  orderBroken' = orderBroken \\
  waitHeld' = waitHeld \\
  wakeHeld' = wakeHeld
\end{schema}

\begin{schema}{AcquireUnordered}
  \Delta LS \\
  t? : T
\where
  ts~t? = trun \\
  act~(pr~t?, pc~t?) = grab \\
  lck~(pr~t?, pc~t?) \notin \dom owner \\
  \exists l : \dom owner | owner~l = t? \spot
    lvl~l \geq lvl~(lck~(pr~t?, pc~t?)) \\
  owner' = owner \cup \{ lck~(pr~t?, pc~t?) \mapsto t? \} \\
  pc' = pc \oplus \{ t? \mapsto pc~t? + 1 \} \\
  ts' = ts \\
  pr' = pr \\
  woken' = woken \\
  orderBroken' = fset \\
  waitHeld' = waitHeld \\
  wakeHeld' = wakeHeld
\end{schema}

\begin{zed}
  Acquire \defs AcquireOrdered \lor AcquireUnordered
\end{zed}

Acquire is guarded by ``the lock is free'': mutual exclusion, assumed as a
lock is modelled and not checked.  The two disjuncts differ only in the flag:
an acquisition made while the thread holds a lock of equal or higher rank
sets \texttt{orderBroken}, and both proceed, so that a broken order is
reported by the goal instead of hidden by a disabled step.

\begin{schema}{Release}
  \Delta LS \\
  t? : T
\where
  ts~t? = trun \\
  act~(pr~t?, pc~t?) = drop \\
  owner' = \{ lck~(pr~t?, pc~t?) \} \ndres owner \\
  pc' = pc \oplus \{ t? \mapsto pc~t? + 1 \} \\
  ts' = ts \\
  pr' = pr \\
  woken' = woken \\
  orderBroken' = orderBroken \\
  waitHeld' = waitHeld \\
  wakeHeld' = wakeHeld
\end{schema}

\begin{schema}{WaitFree}
  \Delta LS \\
  t? : T
\where
  ts~t? = trun \\
  act~(pr~t?, pc~t?) = wt \\
  t? \in woken \\
  t? \notin \ran owner \\
  pc' = pc \oplus \{ t? \mapsto pc~t? + 1 \} \\
  ts' = ts \\
  pr' = pr \\
  owner' = owner \\
  woken' = woken \setminus \{ t? \} \\
  orderBroken' = orderBroken \\
  waitHeld' = waitHeld \\
  wakeHeld' = wakeHeld
\end{schema}

\begin{schema}{WaitHeld}
  \Delta LS \\
  t? : T
\where
  ts~t? = trun \\
  act~(pr~t?, pc~t?) = wt \\
  t? \in woken \\
  t? \in \ran owner \\
  pc' = pc \oplus \{ t? \mapsto pc~t? + 1 \} \\
  ts' = ts \\
  pr' = pr \\
  owner' = owner \\
  woken' = woken \setminus \{ t? \} \\
  orderBroken' = orderBroken \\
  waitHeld' = fset \\
  wakeHeld' = wakeHeld
\end{schema}

\begin{zed}
  Wait \defs WaitFree \lor WaitHeld
\end{zed}

The caller's wait ends once its own event is set, and consumes that event
alone: another caller's signal does not end it.  A wait made while the thread
holds any lock sets \texttt{waitHeld}: C6 says no lock is held while waiting.

\begin{schema}{Deadline}
  \Delta LS \\
  d? : T
\where
  ts~d? = trun \\
  act~(pr~d?, pc~d?) = wt \\
  d? \notin woken \\
  woken' = woken \cup \{ d? \} \\
  ts' = ts \\
  pr' = pr \\
  pc' = pc \\
  owner' = owner \\
  orderBroken' = orderBroken \\
  waitHeld' = waitHeld \\
  wakeHeld' = wakeHeld
\end{schema}

The deadline of one caller fires.  It is enabled for a thread at its wait
whose event is not yet set, so no wait outlasts its own deadline; it sets that
thread's event alone and takes no lock.

\begin{schema}{SignalFree}
  \Delta LS \\
  t? : T \\
  w? : T
\where
  ts~t? = trun \\
  act~(pr~t?, pc~t?) = sg \\
  (tbl , t?) \notin owner \\
  ts~w? = trun \\
  woken' = woken \cup \{ w? \} \\
  pc' = pc \oplus \{ t? \mapsto pc~t? + 1 \} \\
  ts' = ts \\
  pr' = pr \\
  owner' = owner \\
  orderBroken' = orderBroken \\
  waitHeld' = waitHeld \\
  wakeHeld' = wakeHeld
\end{schema}

\begin{schema}{SignalHeld}
  \Delta LS \\
  t? : T \\
  w? : T
\where
  ts~t? = trun \\
  act~(pr~t?, pc~t?) = sg \\
  (tbl , t?) \in owner \\
  ts~w? = trun \\
  woken' = woken \cup \{ w? \} \\
  pc' = pc \oplus \{ t? \mapsto pc~t? + 1 \} \\
  ts' = ts \\
  pr' = pr \\
  owner' = owner \\
  orderBroken' = orderBroken \\
  waitHeld' = waitHeld \\
  wakeHeld' = fset
\end{schema}

\begin{zed}
  Signal \defs SignalFree \lor SignalHeld
\end{zed}

A resolution wakes one waiter, the owner of the call it resolved: \texttt{w?}
is any running thread, so the waiter may be anywhere in its program, or
the resolved call's caller may already be gone.  A wake made while the thread holds the table
lock sets \texttt{wakeHeld}: C6 says the resolution wakes the waiter after the
table lock is released.

\begin{schema}{Finish}
  \Delta LS \\
  t? : T
\where
  ts~t? = trun \\
  act~(pr~t?, pc~t?) = fin \\
  ts' = ts \oplus \{ t? \mapsto tdone \} \\
  pc' = pc \\
  owner' = owner \\
  pr' = pr \\
  woken' = woken \\
  orderBroken' = orderBroken \\
  waitHeld' = waitHeld \\
  wakeHeld' = wakeHeld
\end{schema}

\begin{schema}{Halt}
  \Xi LS
\where
  \forall t : T \spot ts~t = tdone
\end{schema}

Halt is enabled only when every thread has finished, so a state in which some
thread is not done and no thread can move is a deadlock of the model, and
probcli reports it.  A thread holding the table lock waits for nothing; a
thread holding the registry lock waits only for the table lock; a thread
holding the store lock waits only for the registry and table locks; a caller
waits holding nothing, and its wait ends at the deadline.  The acquisition
graph is acyclic.

\begin{verbatim}
  fuzz docs/client_calls_locks.tex
  scripts/verify_client_calls_model.sh
\end{verbatim}

\end{document}
